\section{What Is a Runic Warrior?}

A Runekeeper is defined by two things: a Codex that records their Rites, and a Thiasos that binds them to their Patron. Most Runekeepers carry their Codex as a book—vellum pages sewn between leather boards, ink and blood pressed into the margin. Most keep a Thiasos that chitters or growls or watches from the shadow of a hood.

But stone does not bind well to parchment. Steel does not read. The wild does not keep a library.

A Runic Warrior writes their Rites on a different surface. A gorget, engraved with vows. A living staff, bark split by carved runes. The skin itself, needle and ash threading the flesh. Their Codex is not read by candlelight in a scriptorium. It is felt. It is worn. It is known in the same way a warrior knows the heft of a sword.

And their Thiasos?

A warhorse that has carried them through three campaigns, loyal not out of magic but out of trust. A spirit-hound that runs beside them, no leash, only kinship. A hollowed bird with milk-white eyes, silent witness to every wound and every oath. Or, stranger still, a sword that speaks in the voice of a dead master—never acting on its own, but always \textit{remembering}.

This is not a path for scholars. It is not a path for priests. It is a path for those who answer a question with steel, who swear an oath with blood, who carry their Patron's will in the weight of their shoulders and the scars on their hands.

\section{The Philosophy of Physical Obligation}

Every Runekeeper owes. That is the first truth of the covenant. But where a traditional Runekeeper owes in ink—pages turned, Rites recorded, witnesses notarized—a Runic Warrior owes in \textit{flesh}.

A scar is a debt paid. A limp is a Rite that cost more than expected. A callus is a promise kept.

This is not a metaphor. When a Runic Warrior invokes a Rite, they do not reach for a book. They reach for a scar. They touch the engraving on their gauntlet. They tap the rune carved into their staff. The Codex is not separate from them. It is a part of them.

The price of this intimacy is vulnerability. A book can be replaced. A scar cannot. A shield can be reforged. A hand, once lost, is lost.

But the strength of it is the same. A Runic Warrior's power cannot be taken by a thief in the night. It cannot be burned by a censer. It lives in the body, and the body does not forget.

\section{How to Use This Expansion}

\subsection{For Players}

You are building a warrior who is also a Runekeeper. You are not a bookish scholar. You are not a courtier. You are the one who stands at the front, who holds the line, who charges when others hesitate.

This expansion provides:

\begin{itemize}
\item New talents that emphasize physical action and resilience.
\item New Rites designed for combat and threshold work.
\item Guidelines for creating a Codex that is a weapon, a shield, or a part of your body.
\item Rules for a Thiasos that fights beside you, whether it has flesh or a blade.
\end{itemize}

You do not need to abandon the core Runekeeper rules. This expansion builds on them. You will still track Obligation. You will still invoke Rites. But your tools are different, and your costs are written on your skin.

\subsection{For Game Masters}

You are incorporating a new kind of Runekeeper into your world. These are not wizards in armor. They are warriors who have made a covenant with forces that demand physical proof of loyalty.

Consider:

\begin{itemize}
\item How do the Chain-Lanterns view a Runekeeper whose Codex is a sword? Does that make them harder to prosecute, or more suspicious?
\item How do the Daughters of the Cord treat a Bloodsworn who wears their Rites as tattoos? Is that kinship or competition?
\item How does a Lethai Ghe'hai operative navigate courts where silk masks are expected and open steel is an insult?
\end{itemize}

This expansion also introduces optional rules for sentient Thiasos—weapons that speak, shields that remember. Use them as allies, as complications, or as narrative devices that carry secrets and opinions.

\subsection{Cultural Integration}

Runic Warriors appear in every region of Fate's Edge, but their forms differ.

\begin{itemize}
\item In \textbf{Aeler}, a Runic Warrior might be a Cairn-Knight, Rites carved into a burial stone, their Thiasos an ancestral shade that fights beside them.
\item In \textbf{Sekogo}, they are Beast Warriors, wearing the skin of a predator, their Codex a tooth, their Thiasos the spirit of the hunt.
\item In \textbf{Ashaan}, they are Shadow Warriors, moving unseen, their Codex a black feather, their Thiasos a raven that carries whispers.
\item In \textbf{Lethai}, they are Ghe'hai operatives—silks, steel, and silence. Their Codex is a woven mask. Their Thiasos is a spider-spirit that weaves threads between enemies.
\end{itemize}

This expansion does not favor one region over another. It provides the tools to build a Runic Warrior from any people, under any Patron.

The thread held. The water carried. That is not a Rite. That is grace.
